\documentclass[a4paper]{article}
\usepackage{amsfonts}
\usepackage{amsmath}
\usepackage{amssymb}
\usepackage{graphicx}  

\newcommand{\degree}{^{\circ}}
\newcommand{\cis}{\operatorname{cis}}
\newcommand{\C}{\mathbb{C}}

\begin{document}
  
\noindent \large Auteur: Rune Suy en Torben Petré \\
\noindent \large Studierichting: Informatica\\
\noindent \large Oefeningenreeks 1D nr. 20.6\\

\medskip

\normalsize

Bepaal de nulpunten en de ontbinding in factoren in $\mathbb{R}[z]$ van: \\

$F(x) = z^8 + z^4 + 1$\\

Manier 1:\\

stel $t = z^4$\\

dan $F(x) = t^2 + t + 1$\\

$D = 1- 4 = -3$\\

$t_1 =  \dfrac{-1 - \sqrt{-3}}{2} = \dfrac{-1 - i\sqrt{3}}{2} = \cis\left(\dfrac{-2\pi}{3}\right)$\\

$t_2 =  \dfrac{-1 + \sqrt{-3}}{2} = \dfrac{-1 + i\sqrt{3}}{2} = \cis\left(\dfrac{2\pi}{3}\right)$\\

$z_1 = \sqrt[4]{t_1} = \sqrt[4]{1} \cis\left(\dfrac{\dfrac{-2\pi}{3} + k2\pi}{4}\right)$ met $k \in \{0,1,2,3\}$\\

$(z_1)_1 = \sqrt[4]{1} \cis\left(\dfrac{\dfrac{-2\pi}{3}}{4}\right) = \cis\left(\dfrac{-\pi}{6}\right) = \dfrac{\sqrt{3}}{2} - \dfrac{1}{2}i$\\

$\Rightarrow (z_2)_1 = \dfrac{\sqrt{3}}{2} + \dfrac{1}{2}i$\\

$(z_1)_2 = \sqrt[4]{1} \cis\left(\dfrac{\dfrac{-2\pi}{3} + 2\pi}{4}\right) = \cis\left(\dfrac{\pi}{3}\right) = \dfrac{1}{2} + \dfrac{\sqrt{3}}{2}i$\\

$\Rightarrow (z_2)_2 = \dfrac{1}{2} - \dfrac{\sqrt{3}}{2}i$\\

$(z_1)_3 = \sqrt[4]{1} \cis\left(\dfrac{\dfrac{-2\pi}{3} + 4\pi}{4}\right) = \cis\left(\dfrac{5\pi}{6}\right) = \dfrac{-\sqrt{3}}{2} + \dfrac{1}{2}i$\\

$\Rightarrow (z_2)_3 = \dfrac{-\sqrt{3}}{2} - \dfrac{1}{2}i$\\

$(z_1)_4 = \sqrt[4]{1} \cis\left(\dfrac{\dfrac{-2\pi}{3} + 6\pi}{4}\right) = \cis\left(\dfrac{4\pi}{3}\right) = \dfrac{-1}{2} - \dfrac{\sqrt{3}}{2}i$\\

$\Rightarrow (z_2)_4 = \dfrac{-1}{2} + \dfrac{\sqrt{3}}{2}i$\\

Er bestaan geen reële nulpunten. Als we elk nulpunt uit $z_1$ vermenigvuldigen \\

met zijn complementaire $z_2$ bekomen we de ontbinding in factoren in $\mathbb{R}[z]$.\\

$F(x) = \left[ z - \left(\dfrac{\sqrt{3}}{2} - \dfrac{1}{2}i \right)\right] \cdot \left[z - \left(\dfrac{\sqrt{3}}{2} + \dfrac{1}{2}i\right)\right] \cdot \left[z- \left(\dfrac{1}{2} + \dfrac{\sqrt{3}}{2}i\right)\right] \cdot \left[z- \left(\dfrac{1}{2} - \dfrac{\sqrt{3}}{2}i\right)\right]$\\

$\cdot \left[z- \left(\dfrac{-\sqrt{3}}{2} + \dfrac{1}{2}i\right)\right] \cdot \left[ z - \left(\dfrac{-\sqrt{3}}{2} - \dfrac{1}{2}i\right)\right] \cdot \left[z - \left(\dfrac{-1}{2} - \dfrac{\sqrt{3}}{2}i\right)\right] \cdot \left[z - \left(\dfrac{-1}{2} + \dfrac{\sqrt{3}}{2}i\right)\right]$\\

$= (z^2 - \sqrt{3}z + 1) \cdot (z^2 - z + 1) \cdot (z^2 + \sqrt{3}z + 1) \cdot (z^2 + z + 1) $\\

\bigskip

Andere methode:\\

$F(x) = z^8 + z^4 + 1$\\

= $z^8 + 2z^4 - z^4 + 1$\\

= $(z^4 + 1)^2 - z^4$\\

= $[(z^4 + 1) - z^2] \cdot [(z^4 + 1) + z^2]$\\

= $[z^4 + 1 + 2z^2 - 3z^2] \cdot [z^4 + 1 + 2z^2 - z^2]$\\

= $[(z^2 + 1)^2 - 3z^2] \cdot [(z^2 + 1)^2 - z^2]$\\

= $[(z^2 + 1) - \sqrt{3}z] \cdot [(z^2 + 1) + \sqrt{3}z] \cdot [(z^2 + 1) - z] \cdot [(z^2 + 1) + z]$\\

= $(z^2 - \sqrt{3}z + 1) \cdot (z^2 + \sqrt{3}z + 1) \cdot (z^2 - z + 1) \cdot (z^2 + z + 1)$\\

Ook hieruit kun je afleiden dat er geen reële nulpunten zijn. Je kunt wel 8\\

nulpunten vinden in $\C$. Dit doen we met de discrimant van de vier\\

resulterende $2^{de}$ graads functies.\\

$(z^2 - \sqrt{3}z + 1)$\\

$\Delta_1 = (-\sqrt{3})^2 - 4 \cdot 1 \cdot 1 = -1$ \\

$\Rightarrow (z_1)_1 = \dfrac{- (-\sqrt{3}) - \sqrt{-1}}{2 \cdot 1} = \dfrac{\sqrt{3}}{2} - \dfrac{1}{2}i$ \\

$\Rightarrow (z_1)_2 = \dfrac{- (-\sqrt{3}) + \sqrt{-1}}{2 \cdot 1} = \dfrac{\sqrt{3}}{2} + \dfrac{1}{2}i$ \\

$(z^2 + \sqrt{3}z + 1)$ \\

$\Delta_2 = (\sqrt{3})^2 - 4 \cdot 1 \cdot 1 = -1$ \\

$\Rightarrow (z_2)_1 = \dfrac{-\sqrt{3} - \sqrt{-1}}{2 \cdot 1} = \dfrac{-\sqrt{3}}{2} - \dfrac{1}{2}i$ \\

$\Rightarrow (z_2)_2 = \dfrac{-\sqrt{3} + \sqrt{-1}}{2 \cdot 1} = \dfrac{-\sqrt{3}}{2} + \dfrac{1}{2}i$ \\

$(z^2 - z + 1)$\\

$\Delta_3 = (-1)^2 - 4 \cdot 1 \cdot 1 = -3$ \\

$\Rightarrow (z_3)_1 = \dfrac{- (-1) - \sqrt{-3}}{2 \cdot 1} = \dfrac{1}{2} - \dfrac{\sqrt{3}}{2}i$ \\

$\Rightarrow (z_3)_2 = \dfrac{- (-1) + \sqrt{-3}}{2 \cdot 1} = \dfrac{1}{2} + \dfrac{\sqrt{3}}{2}i$ \\

$(z^2 + z + 1)$\\

$\Delta_4 = 1^2 - 4 \cdot 1 \cdot 1 = -3$ \\

$\Rightarrow (z_4)_1 = \dfrac{- 1 - \sqrt{-3}}{2 \cdot 1} = \dfrac{-1}{2} - \dfrac{\sqrt{3}}{2}i$ \\

$\Rightarrow (z_4)_2 = \dfrac{- 1 + \sqrt{-3}}{2 \cdot 1} = \dfrac{-1}{2} + \dfrac{\sqrt{3}}{2}i$ \\

\begin{thebibliography}{999}
%\bibitem{a} Referentie
\end{thebibliography}
\end{document} 
